\documentclass{article}

\begin{document}
    \begin{center}
        \Large \textbf{Challenge 4: Unter der Haube von Linux} \\ [6pt]
        \normalsize
        Autor: Maximilian Jakob Maag B.Sc. \\
        Version: 1.0
    \end{center}

    \section{Intro}
    In dieser Challenge beschäftigen Sie sich mit den internen Abläufen von Linux.
    Sie lernen, wie das System startet und wo wichtige Systeminformationen gespeichert sind. \\\\
    Ziel ist es, ein besseres Verständnis für die Funktionsweise eines Linux-Systems zu entwickeln. \\\\

    \section{Lernziele}
    Sie verstehen grundlegende Abläufe innerhalb eines Linux-Systems.
    Verschaffen Sie sich einen Überblick über den Bootprozess und Systemlogs.

    \section{Aufgaben}
    Diese Aufgaben sollen Ihnen dabei helfen, die internen Mechanismen von Linux besser zu verstehen.

    \subsection{Aufgabe 1}
    Was passiert beim Start eines Linux-Systems?

    \subsection{Aufgabe 2}
    Informieren Sie sich über Bootloader und deren Aufgabe.

    \subsection{Aufgabe 3}
    Kompilieren Sie ein einfaches Programm aus dem Quellcode (z.B. ein kleines C-Programm).

    \subsection{Aufgabe 4}
    Untersuchen Sie Log-Dateien in /var/log oder mit journalctl.
    Welche Informationen können dort gefunden werden?

    \subsection{Aufgabe 5}
    Was sind Umgebungsvariablen?
    Zeigen Sie, wie man eine Variable setzt und ausliest.

\end{document}